\documentclass[12pt]{article}
\usepackage[a4paper,margin=1in]{geometry}
\usepackage{amsmath,amssymb}
\usepackage{enumitem}
\usepackage{mathtools}
\usepackage{hyperref}
\usepackage{titlesec}
\usepackage{amsthm}

\title{Review of Egor Jakimov}
\author{Simon Gustafsson}
\date{6th of November 2026}


\begin{document}
\maketitle

The report presents the maximal inequality of Fefferman and
Stein in a careful and rigorous way, and the proofs are generally easy to
follow. 

The duality between $\exp L$ and $L\log^+ L$ after Remark 5 is an important
ingredient in the proof of (iii) and perhaps would fit better as a separate
proposition. At the moment Chebyshev's inequality gets its own proposition
while this is only listed as a fact.

Lemma 6 assumes $f$ positive but is later applied to arbitrary $f_k$. Since
$f_k^* = |f_k|^*$, it would be enough to assume $f_k \geq 0$ at the start of
the proof of Theorem 7, or simply to drop the positivity assumption in Lemma 6.

I see why the order of the proof of Theorem 7 is necessary, since the case $p = r$ is
needed for (ii), which in turn gives the remaining cases of (i). It would be
easier to follow this non-linear structure with clearer markers and boundaries 
between the cases. For example, ``This shows that (i) holds'' on page 6 refers 
only to the case $p = r$ while I understood it to first claim proof of all cases of (i).

Moreover, vague references to the course and to Canvas could be replaced by
references to the course literature, which would make the report more
self-contained. The introduction could also say a few words about the motivation for these results. 
Finally, my opinion is that a short abstract and suitable title could add to the reading experience.

\end{document}